\scanpage{113}
Je dis que cette formule implique que le diagramme similaire au précédent
\diagram{\begin{tikzcd}[column sep=large,row sep=large]
A_1\otimes_k A_2 \arrow[r] \arrow[d,"{\deleted{\unsure{N}}n_1\otimes_k\mathrm{id}_{A_2}}"',"=N_{12}"] & A \arrow[d,no head,"N_2"] \\
A_2 \arrow[r,Rightarrow] & A_2
\end{tikzcd}}
\margin{déduite de la forme $n_1$ par chgt de base $k\mathbin{-}A_2$}
C'est en tous cas vrai pour les éléments de \deleted{la forme} $A_1\otimes_k A_2$ de la forme $x_1\otimes x_2$, car les deux mbs sont alors
\[
 t_1(x_1)\cdot\deleted{\unsure{\alpha_2}}(\unsure{x_2})^{\unsure{2}}
\]
par quadraticité de $N_{1,2},N_2$. Il reste à vérifier qu'on a commutativité pour les \deleted{formes} bil. associées, or pour $N_{12}$ c'est la forme déduite de
\[
 \deleted{(x,y)\longrightarrow}\quad\varphi_{n_1}(x,y)=t_1(x\bar y)\qquad\text{sur }A_1\times A_1,
\]
par chgt. de base $k\longrightarrow A_1$, or on a
\margin{si $x'=\alpha_1(x),\ y'=\alpha_1(y)$}
\[
 \left\{\begin{aligned}
 \deleted{\unsure{x\bar y}}\alpha_1(\bar y)&=\sigma_2(\alpha_1(y))\\
 \underset{\text{dans }A_2}{t_1(x\bar y)}&=T_2(\unsure{\alpha_1}(x\bar y))
 \end{aligned}\right.
\]
donc
\[
 \begin{aligned}
 \varphi_{\unsure{n|n_1}}(x,y)&=T_2(x'\sigma_2(y'))\\
 &\mathrel{\text{''}}\varphi_{N_2}(x',y')
 \end{aligned}\qquad\text{O.K.}
\]
\mspage{112}
